\documentclass[12pt]{article}

\usepackage[margin=1in]{geometry}
\usepackage[T1]{fontenc}
\usepackage[utf8]{inputenc}
\usepackage{setspace}
\usepackage{parskip}

\setlength{\parindent}{0pt}
\onehalfspacing

\begin{document}

\title{America's Path Forward}
\author{}
\date{}
\maketitle

Inhales of the crisp, maple-scented air take me down the gravel footpaths of Minuteman National Park, MA, a small but historical park not far from my hometown. North Bridge arches over a little stream meandering its way through trees with gold-tinted leaves just as it did 250 years ago. The stories of fallen heroes will never be forgotten, and the tranquil depths of the New England forests will never change.

Places of true, eternal serenity like this one can be found throughout our nation. Yet, chaos rages across the Internet, social media, and news groups. Hateful violence, extreme rhetoric, and opposing ideas competing to influence legislation have blurred our vision. This chaos is causing many people to doubt whether or not the nation is heading down the right direction. We indeed have come a long way---250 years---since our founding. So, how can we be able to confidently take steps into our future?

To this, I like to imagine that Uncle Sam has two legs: one called Law and the other called Morals. The only way for the nation to move forward is to walk on both of these legs. Law can uphold order to a certain extent, but beyond the domain of Law, the only way to promote harmony is to rely on Morals. Nowadays, too often do we push for legislation---the Law leg---while leaving the Moral leg to trail behind.

What is ``Moral''? For those of you who like concrete definitions, this subjective word cannot be defined easily. Nevertheless, there is one key driver that pushes our Moral foot forward: civic education.

Education is so fundamental that it is rooted adjacent to our founding values. As the Preamble of the Constitution states, we need to ``secure the Blessings of Liberty to ourselves and our Posterity''. In order to do so, we need to prioritize promoting a civic education that creates respectful, responsible citizens who understand the way U.S. society works. In other words, civic education promotes ``civic participation''---something I first learned in my eighth-grade Civics class.

``How do you engage in civic participation?''

When my Civics teacher asked my class that question, we were unsure of what it meant. He then explained it could be anything from casting a vote on a town election to coaching a youth sports team.

``Civic participation,'' he said, ``means giving to---and therefore investing in---your community.''

Hearing this, I reflected upon myself. Even though I cannot vote yet, I do submit articles to my town's newsletter, perform music at cultural events, attend town meetings, and plant flowers along my neighborhood driveway. Maybe America's 250th birthday means the 250th year of individuals planting seeds of hope for those around them.

At the end of the day, no matter how many laws we tell the government to pass or how heated online debates become, our nation can only function if we citizens engage in civic participation. It is almost like an invisible contract that cannot be written down formally on a piece of legislation nor be enforced through punishments: Individuals get their freedom, and in turn they are expected to bring positivity to society. You don't have to have a ton of money to donate. All you need to do is start small and add something you think would taste good to the soup.

Civic education of our youth will help maintain a civilized and respectful American culture. It will also produce well-rounded citizens that will engage in civic participation. As a middle schooler, I know that too often kids my age live and breathe in our own individual worlds and rarely see ourselves as part of a bigger picture. In addition, due to a lack of civic education, only 1 in 3 of American-born citizens can pass the U.S. Naturalization Test as of 2018 (Woodrow Wilson Foundation), compared to the 95.7\% pass rate of all immigrant applicants in 2022 (U.S. Citizenship and Immigration Services).

In response to this problem, my home state of Massachusetts mandated all 8th graders to take one year of Civics (Massachusetts General Court). Civics class has had one of the biggest impacts on me, because everything I learn in it is applicable everywhere in my daily life. As the year progressed, more and more I realized that I am part of a larger community---my town, county, state, region, and nation. American children in all schools across the nation deserve to receive a Civics education like me.

As the nation with the longest-running written constitution still in operation, the United States has succeeded to sustain its founding values for a very long time. Throughout these past 250 years, our founding ideals of ``life, liberty, and the pursuit of happiness'' have remained steadfast. To ensure they never will change even in the face of future challenges, our founders equipped us with the right to add amendments to the Constitution---a method to move our nation's Law leg forward. Given this, we also have the responsibility to be educated, virtuous citizens and propel Morals leg forward. For me, America's founding ideals mean that anyone has the freedom to choose. Anyone can choose to give something, no matter how small, to society. The gravel path I'm walking on was created for me and all future generations by 250 years of past Americans treading upon it. America's path is one that needs to be walked out by all of us.

\section*{Works Cited}

Massachusetts General Court. ``Chapter 296: An Act to Promote and Enhance Civic Engagement.'' \textit{Massachusetts Legislature (malegislature.gov)}, 8 Nov.\ 2018. \texttt{malegislature.gov/Laws/SessionLaws/Acts/2018/Chapter296}. Accessed 12 Mar.\ 2026.

U.S.\ Citizenship and Immigration Services. \textit{Naturalization Test Performance: National Pass Rate}. 5 Oct.\ 2023. \texttt{www.uscis.gov/citizenship-resource-center/naturalization-related-data-and-statistics/naturalization-test-performance}. Accessed 12 Mar.\ 2026.

Woodrow Wilson Foundation. ``National Survey Finds Just 1 in 3 Americans Would Pass Citizenship Test.'' \textit{PR Newswire}, 3 Oct.\ 2018. \texttt{www.prnewswire.com/news-releases/national-survey-finds-just-1-in-3-americans-would-pass-citizenship-test-300723440.html}. Accessed 12 Mar.\ 2026.

\end{document}